\scanpage{25}
\mspage{- IV-\handwritten{950} -}
$f_{\unsure{o|0}}^{(\handwritten{\lambda})}:S_o^{(\handwritten{\lambda})}\to {S'}_o^{(\handwritten{\lambda})}$ tels que $S'=S'^{(\handwritten{\lambda})}\times_{S^{(\handwritten{\lambda})}}S$ , $g=g^{(\handwritten{\lambda})}\times 1$ et $f_o=f_o^{(\handwritten{\lambda})}\times 1$ . \deleted{\illegible{1}}
Soient $s$ le point fermé de $S$ , $x=f_o(s)\deleted{\unsure{{}_s}}$ , $x_{\handwritten{\lambda}}$ la projection de $x$ dans $S'^{(\handwritten{\lambda})}$ ;
comme $x_{\handwritten{\lambda}}$ est au-dessus du point fermé de $S^{(\handwritten{\lambda})}$ , l'anneau local \inserted{\handwritten{$C_\lambda$}} de $S'^{(\handwritten{\lambda})}$ au
point $x_{\handwritten{\lambda}}$ est une $A_{\handwritten{\lambda}}$-algèbre essentiellement étale ; en outre , comme $f_o^{(\handwritten{\lambda})}(s_{\handwritten{\lambda}})=$
$=x_{\handwritten{\lambda}}$ , on a $\handwritten{\kappa}(x_{\handwritten{\lambda}})=\handwritten{\kappa}(s_{\handwritten{\lambda}})$ , autrement dit , $\unsure{C|S|\mathcal{O}}_{\handwritten{\lambda}}$ est une $A_{\handwritten{\lambda}}$-algèbre locale strictement
essentiellement étale , et par suite (18.6.1) est $A_{\handwritten{\lambda}}$-isomorphe à une
$A$-algèbre $A_{\handwritten{\mu}}$ avec $\handwritten{\mu\geq\lambda}$ . Il y a \deleted{\illegible{2}} \inserted{\handwritten{donc}} un $A_{\handwritten{\lambda}}$-homomorphisme \deleted{\illegible{1}} \inserted{\handwritten{$C_\lambda\to{}^{\mathrm h}A$}} , c'est-à-dire
un $S^{(\handwritten{\lambda})}$-morphisme $h:S\to S'^{(\handwritten{\lambda})}$ tel que $h(s)=x_{\handwritten{\lambda}}$ , et par suite il existe
bien une $S$-section $f$ de $S'$ telle que $f(s)=x$ , ce qui achève de prouver (i) \unsure{;|.}

Pour prouver (ii) , il suffit de noter que l'on a $\operatorname{Hom.loc}({}^{\mathrm h}A,B)\deleted{\unsure{=}}$
\margin{\handwritten{$\simeq$}}
$=\varprojlim\operatorname{Hom.loc}(A_{\handwritten{\lambda}},B)$ et que par (18.6.2) les homomorphismes canoniques
$\operatorname{Hom.loc}(A_{\handwritten{\lambda}},B)\longleftarrow\operatorname{Hom.loc}(A,B)$ sont bijectifs .

\deleted{\illegible{1}} Pour démontrer \deleted{\illegible{1}} (iv) , \deleted{\illegible{2}} notons que pour tout $\handwritten{\lambda}$ et tout
entier $n>0$ , $\underline m_{\handwritten{\lambda}}^n=\underline m^n A_{\handwritten{\lambda}}$ , et $(\underline m.{}^{\mathrm h}A)^n=\underline m^n.{}^{\mathrm h}A=\varinjlim\underline m_{\handwritten{\lambda}}^n$ . On en conclut , par l'exactitude
du foncteur $\varinjlim$ dans la catégorie des $A$-modules , que $ {}^{\mathrm h}A/(\underline m.{}^{\mathrm h}A)^n=$
$=\varinjlim(A_{\handwritten{\lambda}}/\underline m_{\handwritten{\lambda}}^n)$ , et par suite , il suffit de montrer que , pour tout entier $n$,
et tout indice $\handwritten{\lambda}$ , l'homomorphisme $A/\underline m^n\to A_{\handwritten{\lambda}}/\underline m_{\handwritten{\lambda}}^n$ est bijectif . Or , \deleted{\illegible{1}} cela
est vrai par hypothèse pour $n=1$ ; d'autre part , puisque $A_{\handwritten{\lambda}}$ est un $A$-module
plat , on a
\[
\underline m_{\handwritten{\lambda}}^n/\underline m_{\handwritten{\lambda}}^{n+1}
=(\underline m^n/\underline m^{n+1})\otimes_A A_{\handwritten{\lambda}}
=(\underline m^n/\underline m^{n+1})\otimes_{A/\underline m}(A_{\handwritten{\lambda}}/\underline m A_{\handwritten{\lambda}})
\]
et comme $A/\underline m\to A_{\handwritten{\lambda}}/\underline m A_{\handwritten{\lambda}}=A_{\handwritten{\lambda}}/\underline m_{\handwritten{\lambda}}$ est bijectif , l'homomorphisme $\underline m^n/\underline m^{n+1}\to\underline m_{\handwritten{\lambda}}^n/\underline m_{\handwritten{\lambda}}^{n+1}$
est lui aussi bijectif ; la conclusion résulte donc de Bourbaki,\uline{Alg.comm.},
chap.III,§ 2,n°8,\handwritten{cor. 3 du th. 1} .

Enfin , si $A$ est hensélien \deleted{[ et si $A'=B_{\underline n}$ est une $A$-algèbre locale
strictement essentiellement étale , où $B$ est une $A$-algèbre étale , il
résulte de \deleted{\illegible{1}} (18.6.2) qu'il existe une \deleted{\illegible{1}} $\operatorname{Spec}(A)$-\uline{section de}
$\operatorname{Spec}(B)$ prenant la valeur $\underline n$ au point fermé de $\operatorname{Spec}(A)$ ]} ; on conclut
\deleted{alors} \inserted{\handwritten{de la remarque qui \unsure{s|suit}}} de \deleted{(17.4.1)} \inserted{\handwritten{18.6.2}} que \deleted{\illegible{2}} \deleted{l'}\inserted{\handwritten{les}} homomorphisme\handwritten{s} \deleted{\illegible{1}} $A\to A\deleted{'}_{\handwritten{\lambda}}$ \deleted{est} \inserted{\handwritten{sont}}
bijectif\handwritten{s}, ce qui prouve (vi) par définition de $ {}^{\mathrm h}A$ .

(18.6.7) Soient maintenant $A$ un anneau \uline{semi-local} , $\underline m_i$ $(1\leq i\leq r)$ ses idéaux
maximaux . On appelle \uline{hensélisé} de $A$ et l'on note encore $ {}^{\mathrm h}A$ l'anneau \uline{produit}
$\handwritten{\prod_i}{}^{\mathrm h}(A_{\underline m_i})$ des hensélisés des anneaux locaux $A_{\underline m_i}$ . C'est \inserted{\handwritten{un $A$-module fidèlement plat et}} une $A$-algèbre semi-locale
, dont les idéaux maximaux sont les $\underline m_i.{}^{\mathrm h}A$ en vertu de (18.\unsure{6|5}.6,(iii));
en outre , si $\underline r=\handwritten{\bigcap_i}\underline m_i$ est le radical de $A$ , il résulte de ce qui précède
que $\underline r.{}^{\mathrm h}A$ est le radical de $ {}^{\mathrm h}A$ et que l'application $A/\underline r\to{}^{\mathrm h}A/\underline r.{}^{\mathrm h}A$ est bijective
. Comme le séparé complété $\hat A$ de $A$ pour la topologie $\underline r$-préadique est le
produit des séparés complétés $\hat A_{\underline m_i}$ (Bourbaki,\uline{Alg.comm.},chap.III,§ 2,n°13 \handwritten{prop.
18}) l'homomorphisme canonique $\hat A\to({}^{\mathrm h}A)\handwritten{{}^{\wedge}}$ est bijectif par (18.6.6,(iv)),
% archive sheet number: 24
